% Imporditud: tools/import_eksperimendid.py
% Allikas: Eesti füüsikaolümpiaadi eksperimendi ülesannete
%   kogu (Rael Kalda, 2023), ülesanne Ü129, lahendus L129.
\setAuthor{EFO žürii}
\setRound{lõppvoor}
\setYear{2004}
\setNumber{G E2}
\setDifficulty{4}
\setTopic{elektriahelad}
\setVahendid{taskulambipirn, lapik patarei, takisti, krokodill, tester}
\setPunktid{12}
\setAllikas{Eksperimendi kogu Ü129, lahendus lk 100}
\setLahendusseis{toores}

\prob{Pirn}
Taskulambipirni hõõgniidi takistus on toitepinge U astmefunktsioon R = $R_0$ + cU n, kus c ja n on konkreetset lampi iseloomustavad konstandid. Määrata antud lambi jaoks c ja n.
\textit{Märkus:} ed: Eeldada, et voltmeetrina kasutatava testri takistus on 1 M$\Omega$; kui testrit kasutada ampermeetrina mõõtepiirkonnal 10 A, võib tema takistuse lugeda nulliks. Oommeetrina toimiva testri töövool on 1,2 mA.

\hint

\solu
Ühendame jadamisi patarei, lambi ja testri kui ampermeetri ning registreerime lambi töövoolu $I_1$. Seejärel lisame samasse ahelasse takisti ning mõõdame töövoolu $I_2$. Nüüd teeme samad katsed ilma ampermeetrita, kuid mõõtes testri kui voltmeetriga pinget lambil. Saame vastavad pinge väärtused $U_1$ ja $U_2$. Kuna ampermeetri takistus ja voltmeetri juhtivus (st takistuse pöördväärtus) on tühiselt väike, siis võime eeldada, et voolutugevused on samad, mis eelmistes katsetes. Arvutame Ohmi seadusest lambi takistused mõlemal toitepingel ($R_1$ ja $R_2$). Mõõdame lambi takistuse testri kui oommeetriga. Kuna testri kui oommeetri (ülesandes antud) töövool on sadu kordi väiksem juba meile teada olevast lambi töövoolust, siis saame määrata suuruse $R_0$ (tõsi küll, suure piirveaga, kuid lõpptulemuses kajastub see vähe). Logaritmime seost R $-$ $R_0$ = cU n ja saame ln(R $-$ $R_0$) = ln c + n ln U . Saadud kahe katsepunkti alusel koostame võrrandisüsteemi:

ln($R_1$ $-$ $R_0$) = ln c + n ln $U_1$, ln($R_2$ $-$ $R_0$) = ln c + n ln $U_2$. Olles lahutanud esimesest võrrandist teise, avaldame n ja c kujul

n = ln($R_1$ $-$ $R_0$) $-$ ln($R_2$ $-$ $R_0$) ; c = $R_1$ $-$ $R_0$ . ln $U_1$ $-$ ln $U_2$ Un

Olles teostanud mõõtmised ja arvutused, saame n $\approx$ 0,6 ning c $\approx$ 5,5 $\Omega$ V$^{-0}$,6.
\probend
